Diffraction is a powerful method of analyzing the behaviour of waves
including how they propagate and interact with materials.
For light, common ways of observing this behaviour are single slit,
double slit, and diffraction grating experiments which can be used to
analyze the spectrum of an incident beam.
When a beam enters a medium, the velocity and wavelength of the light
change; the ratio of the speed of light in a vacuum to this velocity in
the medium is called the index of refraction and is unique to a
medium's optical density.
This experiment focuses on using two lasers of known wavelengths to
analyze the refractive indices of a variety of media with theoretical
indices ranging from about 1 to 1.5 at standard atmosphere,
temperature, and pressure.
These materials were placed between a diffraction grating and a screen
to analyze the change in the light's velocity in the medium by
observing the change in the diffraction pattern.
This change in the degrees of diffraction was recorded using a
motorized light-dependent resistor controlled by an Arduino, from which
the refractive indices of air, water, ethanol, and olive oil were
calculated to within 3.2\%, 6.8\%, 10.4\%, and 20.7\% of their expected
values respectively.
These results demonstrate that diffraction-based refractometry provides
an effective and accessible method of determining select optical
properties of non-opaque media.

